\documentclass{article}

\usepackage[utf8]{inputenc}     
\usepackage[T1]{fontenc}
\usepackage[magyar]{babel}

\usepackage{amsmath}
\usepackage{amsthm}
\usepackage{amssymb}

\title{Alkalmazott algebra}
\author{Rovenszky Robin\thanks{Rovenszky Robin fordításával}}
\date{Utoljára frissítve: 2026. március 16.}


\begin{document}

\maketitle

\section{Binomiális tétel}
\begin{theorem}[Binomiális tétel]
A binomiális tétel a binomiális kifejezések hatványainak algebrai kibővítését írja le.
\begin{equation}
    (x+y)^n = \binom{n}{0}x^ny^0 + \binom{n}{1}x^{n-1}y^1 + \dots + \binom{n}{n}x^0y^n.
\end{equation}
\end{theorem}

\section{Feladatgyűjtemény}

\subsection{Algebra 101}

\subsubsection{1.1. feladat}
Bővítsd ki: $(a\pm b)^5$.
\begin{equation}
    (a \pm b)^5
\end{equation}
A binomiális tétel (1. tétel) alkalmazásával:
\begin{equation}
    (a \pm b)^5 = \binom{5}{0}a^5 \pm \binom{5}{1}a^4b + \binom{5}{2}a^3b^2 \pm \binom{5}{3}a^2b^3 + \binom{5}{4}ab^4 \pm \binom{5}{5}b^5.
\end{equation}

\subsubsection{1.2. feladat}
Bővítsd ki: $(x\pm 1)^5$.
\begin{equation}
    (x \pm 1)^5
\end{equation}
A binomiális tétel (1. tétel) alkalmazásával:
\begin{equation}
    (x \pm 1)^5 = \binom{5}{0}x^5\pm \binom{5}{1}x^4 + \binom{5}{2}x^3 \pm \binom{5}{3}x^2 + \binom{5}{4}x \pm \binom{5}{5}.
\end{equation}

\subsubsection{1.3. feladat}
Bővítsd ki: $(a+b+c+d)^2$.
\begin{equation}
    (a+b+c+d)^2 = a^2 + b^2 + c^2 + d^2 + 2ab + 2ac + 2ad + 2bc + 2bd + 2cd.
\end{equation}

\subsubsection{1.4. feladat}
Írd át az $a^2 + ab + a + b$ kifejezést két elsőfokú polinom szorzataként.
\begin{equation}
    a^2 + ab + a + b = a(a+b) + (a+b) = (a + b)(a + 1).
\end{equation}

\subsubsection{1.5. feladat}
Írd át az $x^3 + 2x^2 + 2x + 1$ kifejezést egy másod- és egy elsőfokú polinom szorzataként. (Nem írható át három elsőfokú polinom szorzataként az algebra alaptétele és absztrakt algebra / testelmélet miatt.)

\begin{equation}
    x^3 + 2x^2 + 2x + 1 = (x^3 + x^2 + x) + (x^2 + x + 1) = x(x^2 + x + 1) + (x^2 + x + 1) =
\end{equation}
\begin{equation}
    = (x^2 + x + 1)(x + 1).
\end{equation}

\subsubsection{1.6. feladat}
Írd át az $x^3 + 6x^2 + 12x + 8$ kifejezést polinomok szorzataként.\\
1. módszer:\\
Észrevehetjük, hogy a 1, 6, 12, 8 egy tökéletes köbhöz tartozó binomiális kifejezés mintáját követi. Használjuk az azonosságot:
\begin{equation}
    a^3 + 3a^2b+3ab^2+b^3 = (a+b)^3.
\end{equation}
Tehát:
\begin{equation}
    x^3 + 6x^2 + 12x + 8 = (x+2)^3.
\end{equation}

2. módszer:\\
Észrevesszük, hogy $-2$ a polinom gyöke (azaz a kifejezés $0$-t ad ezen a ponton), így $(x+2)$ tényező. A polinom $(x+2)$-vel való osztása:
\begin{equation}
    x^2 + 4x + 4 = (x+2)^2,
\end{equation}
így
\begin{equation}
    x^3 + 6x^2 + 12x + 8 = (x+2)(x+2)^2 = (x+2)^3.
\end{equation}

\subsubsection{1.7. feladat}
Oldd meg az $ab + 2a - 3b = 14$ egyenletet, ahol $a, b \in \mathbb{N^+}$.

Csoportosítsuk az $a$-t tartalmazó tagokat:
\[
    ab + 2a - 3b = 14
\]
\[
    a(b+2) - 3b = 14
\]
Hozzuk át a $3b$-t a másik oldalra:
\[
    a(b+2) = 3b + 14
\]
Fejezzük ki $a$-t:
\[
    a = \frac{3b+14}{b+2} 
\]
Átalakítva a számlálót:
\[
    3b+14 = 3(b+2) + 8
\]
Így:
\[
    a = 3 + \frac{8}{b+2}.
\]

Ahhoz, hogy $a$ egész legyen, $b+2$-nek osztania kell 8-at.\\
8 osztói: $1, 2, 4, 8$. Mivel $b \ge 1$, $b+2 \ge 3$, így két lehetőség marad:
\begin{itemize}
    \item $b + 2 = 8 \implies b = 6$
    \item $b + 2 = 4 \implies b = 2$
\end{itemize}

Ezeket a $b$ értékeket visszahelyettesítve $a$-ba, a két megoldás: $(5,2)$ és $(4,6)$.

\subsubsection{1.8. feladat}
Határozd meg a $2x^2 + 2y^2 - 2xy + 4x -6y + 8$ kifejezés minimumát.

\subsubsection{1.9. feladat}
Bizonyítsd az alábbi egyenlőtlenséget:
\[
(a+b+c)\left(\frac{1}{a} + \frac{1}{b} + \frac{1}{c}\right) \ge 9
\]
ahol $a, b, c > 0$.

\subsubsection{2.1. feladat}
Bővítsd ki: $(2x-3)^3$.

\subsubsection{2.2. feladat}
Bővítsd ki: $(a-1)(b-1)(c-1)$.

\subsubsection{2.3. feladat}
Írd át a következő kifejezést szorzatként:
\begin{equation}
    (a+b)^2 - (a-b)^2.
\end{equation}

\subsubsection{2.4. feladat}
Írd át a következő kifejezést szorzatként:
\begin{equation}
    x^2y+xy^2-x^2+y^2.
\end{equation}

\subsubsection{2.5. feladat}
Írd át az $a^4 - 1$ kifejezést szorzatként.

Négyzetek különbsége alapján:

\subsubsection{2.6. feladat}
Oldd meg az alábbi egyenletet:
\begin{equation}
    4ab-6a+8b = 14.
\end{equation}

\subsubsection{2.7. feladat}
Határozd meg a következő kifejezés minimumát:
\[
5x^2 + 2y^2 -2x +4y + 4xy + 1.
\]

\subsubsection{2.8. feladat}
Bizonyítsd az alábbi egyenlőtlenséget:
\[
4x+\frac{1}{x} \ge 4
\]
ahol $x > 0$.

\end{document}